\SourceLeaf{415}%
derselben folgende anwenden nach \([6]\)
\[\BareRoot{\left\{rr-\left(\frac{g\mp F}{\lambda}\right)^2\right\}}\]
\par
Wir haben also auch hier zu einer Wurzel zwei Auflösungen, nemlich durch die symmetrischen Lagen der Punkte \((1),(2),(3)\) von \(\mathrm{(II)}\) auf entgegengesetzten Seiten gleich weit abstehend; für \(x=+1\) ist der Sinn der positiven Richtung in \(\mathrm{(I)}\) derselbe wie in \(\mathrm{(II)}\), für \(x=-1\) verkehrt. Nur der einzige Fall, wo ohne Rücksicht auf das Zeichen \(\lambda r=g-F\) (für \(x=+1\)) oder \(=g+F\) (für \(x=-1\)) ist, muss ausgenommen werden, indem dann beide symmetrische Lagen von \(\mathrm{(I)}\) in Eine, nemlich mit der Linie \(\mathrm{(II)}\) selbst, zusammenfallen.

Auszuschliessen sind offenbar von den Wurzeln der Gleichung \([7]\) nicht blos die imaginären, sondern auch die ausserhalb der Grenzen \(-1\) und \(+1\) liegenden reellen, und die Wurzel \(+1\) oder \(-1\) selbst, wenn \(\lambda r\) ohne Rücksicht auf das Zeichen beziehungsweise kleiner ist als \(g-F\) oder \(g+F\). Es lässt sich übrigens beweisen, dass allemal wenigstens eine der drei Wurzeln in die Kategorie der auszuschliessenden gehört, und also überhaupt niemals mehr als vier verschiedene Auflösungen durch reelle Coordinaten statt haben können. Genau genommen, bildet zwar \textit{ein ganz singulärer} Fall in so fern eine Ausnahme des erstern Satzes, als dabei keine Wurzel ausgeschlossen wird. Der singuläre Fall ist nemlich der schon oben erwähnte, wo für \(x=\pm1\) die Ordinate \(=0\) wird, und wo (wie sich leicht beweisen lässt) die betreffende Wurzel zweimal gilt, d.\,i. wo das Glied linkerseits des Gleichheitszeichens in der Gleichung \([7]\) den Factor \((x\mp1)^2\) enthält; die Gleichung hat dann also nur zwei ungleiche Wurzeln, von denen die zweite allerdings auch eine zulässige sein kann. Der Schlussfolge selbst thut demnach dieser Ausnahmefall keinen Eintrag.

Endlich muss noch bemerkt werden, dass auch unter den Auflösungen in reellen Zahlen physisch unzulässige sein können. Es ist nemlich nicht der ganze Kreis, welcher aus dem Mittelpunkte \((1)\) durch \((4)\) und \((5)\) beschrieben ist, der geometrische Ort des Punktes \((1)\), sondern nur der auf der positiven Seite von \((4),(5)\) liegende Bogen, wenn \(\vartheta\) unter \(180^0\), und der auf der negativen Seite liegende, wenn \(\vartheta\) überstumpf ist; dasselbe gilt von den beiden andern Kreisen. Diese physische Bedingung ist aber in unserer Auflösung noch nicht berücksichtigt. Unter den verschiedenen, in reellen Zahlen gefundenen Auflösungen sind also nur diejenigen zulässig, wo die für die Abscissen der drei Punkte \((1),(2),(3)\)
